\documentclass{article}
\usepackage{geometry}
\geometry{margin=1.5cm, vmargin={0pt,1cm}}
\setlength{\topmargin}{-1cm}
\setlength{\headheight}{12.5pt}
\setlength{\paperheight}{29.7cm}
\setlength{\textheight}{25.3cm}

% useful packages.
\usepackage[UTF8]{ctex}
\usepackage{fancyhdr}
\usepackage{layout}
\usepackage{luacode}

% import common macros
% % % % % % % % % % % % % % % % % % % % % % % % % % % % % % % %
% Common Macros for LaTeX
% % % % % % % % % % % % % % % % % % % % % % % % % % % % % % % %

% Useful packages
\usepackage{amsfonts}
\usepackage{amsmath}
\usepackage{amssymb}
\usepackage{amsthm}
\usepackage{enumerate}
\usepackage{graphicx}
\usepackage{multicol}
\usepackage[ruled,linesnumbered]{algorithm2e}

% Math operators and common symbols
\newcommand{\dif}{\mathrm{d}}
\newcommand{\innerProd}[1]{\left\langle #1 \right\rangle}
\newcommand{\difFrac}[2]{\frac{\dif #1}{\dif #2}}
\newcommand{\pdfFrac}[2]{\frac{\partial #1}{\partial #2}}
\newcommand{\T}{\mathrm{T}}
\newcommand{\norm}[1]{\left\| #1 \right\|}
\newcommand{\abs}[1]{\left| #1 \right|}

% Vector and Matrix shortcuts
\newcommand{\vecTwo}[2]{
  \begin{bmatrix}
    #1 \\
    #2
\end{bmatrix}}
\newcommand{\vecThree}[3]{
  \begin{bmatrix}
    #1 \\
    #2 \\
    #3
\end{bmatrix}}
\newcommand{\vecTwoH}[2]{
  \begin{bmatrix}
    #1 & #2
  \end{bmatrix}
}
\newcommand{\vecThreeH}[3]{
  \begin{bmatrix}
    #1 & #2 & #3
  \end{bmatrix}
}

% Common fields
\newcommand{\R}{\mathbb{R}}
\newcommand{\C}{\mathbb{C}}
\newcommand{\Z}{\mathbb{Z}}
\newcommand{\N}{\mathbb{N}}

% Solutions and proofs
\newenvironment{solution}{
\begin{proof}[解]}{
\end{proof}}

% Utility to get environment variables (requires lualatex)
\newcommand{\getEnv}[2]{\directlua{tex.print(os.getenv("#1") or "#2")}}


\newcommand{\tridiag}{\operatorname{tridiag}}
\newcommand{\ii}{\mathrm{i}}

\usepackage{listings}
\usepackage{xcolor}

% Configure listings for Python code highlighting
\lstset{
  language=Python,
  basicstyle=\ttfamily\small,
  keywordstyle=\bfseries\color{blue},
  commentstyle=\color{gray},
  stringstyle=\color{red},
  breaklines=true,
  showstringspaces=false,
  tabsize=4,
  frame=single,
  framesep=5pt,
  rulecolor=\color{gray},
  backgroundcolor=\color{white},
  captionpos=b,
  numberstyle=\tiny\color{gray},
  numbers=left,
  stepnumber=5,
}

\lstnewenvironment{plaintext}[1][]{
  \lstset{
    language=none,
    basicstyle=\ttfamily\small,
    keywordstyle=,
    commentstyle=,
    stringstyle=,
    numbers=none,
    frame=single,
    backgroundcolor=\color{white},
    #1
  }
}{}

% Name and uID
\newcommand{\name}{\getEnv{NAME}{NAME}}
\newcommand{\uid}{\getEnv{UID}{UID}}
\newcommand{\course}{数值代数}
\newcommand{\hwNum}{14}

\begin{document}

\pagestyle{fancy}
\fancyhead{}
\lhead{\name (\uid)}
\chead{\course \#\hwNum}
\rhead{\today}

\section{一般实矩阵的隐式 QR 算法}

求实矩阵的全部特征值及特征向量.

\begin{enumerate}[(1)]
  \item 用熟悉的计算机语言编制利用隐式 QR 算法求一个实矩阵全部特征值和特征向量的通用子程序.
  \item 利用所编制的子程序计算方程
    \[
      x^{41} + x^3 + 1 = 0
    \]
    的全部根.
  \item 设
    \[
      A =
      \begin{bmatrix}
        9.1 & 3.0 & 2.6 & 4.0 \\
        4.2 & 5.3 & 1.7 & 1.6 \\
        3.2 & 1.7 & 9.4 & x \\
        6.1 & 4.9 & 3.5 & 6.2
      \end{bmatrix},
    \]
    试求当 $x = 0.9, 1.0, 1.1$ 时 $A$ 的全部特征值, 并观察特征值的变化情况.
\end{enumerate}

\begin{solution}

  先用 Householder 变换将 $A$ 化为上 Hessenberg 阵，再对其做带 Wilkinson 位移的隐式 QR
  迭代（bulge chasing）.
  收敛后得到复 Schur 分解 $AZ=ZT$，$T$ 的对角元即特征值，特征向量由解上三角系统 $Zy_j=v_j$ 再归一化得到.

  对 $x^{41}+x^3+1=0$，构造对应首一多项式的 $41$ 阶友矩阵，其特征值即多项式的全部根.

  题~(3) 中 $4\times4$ 矩阵的特征值如下：
  \[
    \begin{array}{c|c}
      x & \text{特征值} \\ \hline
      0.9 &
      16.964832670607,\; 7.365629536958,\;
      2.834768896218 \pm 1.007441166353\ii \\
      1.0 &
      16.996516047941,\; 7.343359754830,\;
      2.830062098615 \pm 1.006228905570\ii \\
      1.1 &
      17.028013323148,\; 7.321318091146,\;
      2.825334292853 \pm 1.004978621116\ii
    \end{array}
  \]
  可以看到，当 $x$ 从 $0.9$ 增大到 $1.1$ 时，最大实特征值单调增大；
  另一个实特征值单调减小；复共轭特征值对的实部和虚部绝对值都略有减小，变化较平滑.

  代码与运行结果如下：
  \lstinputlisting[language=Python, caption={一般矩阵隐式 QR 算法及应用}]{exercise.py}
  \lstinputlisting[language={}, caption={一般矩阵隐式 QR 运行结果}]{output.txt}

\end{solution}

\section{实对称矩阵的隐式 QR 算法}

用熟悉的计算机语言编写隐式对称 QR 算法, 求解实对称矩阵的全部特征值和特征向量的通用子程序,
并计算教材第 244 页 1(2) 和 2(2) 的矩阵的全部特征值和特征向量.

第 1(2) 题矩阵为 $n=50,51,\ldots,100$ 阶三对角矩阵
\[
  A_1 = \tridiag(1,4,1)
  =
  \begin{bmatrix}
    4 & 1 \\
    1 & 4 & 1 \\
    & \ddots & \ddots & \ddots \\
    &        & 1 & 4 & 1 \\
    &        &   & 1 & 4
  \end{bmatrix}.
\]
第 2(2) 题矩阵为 $100\times100$ 三对角矩阵
\[
  A_2 = \tridiag(-1,2,-1)
  =
  \begin{bmatrix}
    2 & -1 \\
    -1 & 2 & -1 \\
    & \ddots & \ddots & \ddots \\
    &        & -1 & 2 & -1 \\
    &        &    & -1 & 2
  \end{bmatrix}.
\]

\begin{solution}

  先用 Householder 变换将 $A$ 化为三对角阵 $T=Q_0^\T AQ_0$，再对 $T$ 做带 Wilkinson
  位移的隐式对称 QR 迭代（Givens bulge chasing），累积正交变换即得特征向量.

  常系数三对角 Toeplitz 矩阵 $\tridiag(b,a,b)$ 有精确特征值
  $\lambda_k=a+2b\cos\frac{k\pi}{n+1}$，程序用该公式验证 QR 结果.

  对 $A_1=\tridiag(1,4,1)$，计算了 $n=50,51,\ldots,100$ 全部阶数. $n=100$ 时：
  \[
    \lambda_{\min} = 2.000967435416,\qquad
    \lambda_{\max} = 5.999032564584.
  \]
  对 $A_2=\tridiag(-1,2,-1)$，$n=100$ 时：
  \[
    \lambda_{\min} = 0.000967435416,\qquad
    \lambda_{\max} = 3.999032564584.
  \]
  两个 $100$ 阶例子的最大特征值误差约 $10^{-14}$，特征向量残差约 $10^{-12}$，正交性误差约 $10^{-14}$.

  代码与运行结果如下：
  \lstinputlisting[language=Python, caption={实对称矩阵隐式 QR 算法及应用}]{symmetric_qr.py}
  \lstinputlisting[language={}, caption={实对称矩阵隐式 QR 运行结果}]{symmetric_output.txt}

\end{solution}

\end{document}
